\chapter*{\centering Abstract}
\addcontentsline{toc}{chapter}{Abstract}

\noindent Farmers in the Bengaluru--Kolar belt choose where to sell and whether to wait with only yesterday's price to go on, yet a mandi reports on only part of the days, transport and shrinkage cost money, and a high price at a distant market can be gone when the load arrives. This project built MandiSense AI, a decision-support system that recommends a market and a holding period and attaches a revenue floor tested to hold. It used 63,524 real Agmarknet records (36,821 mandi-days for tomato, onion, potato and ginger). The model was trained on the revenue a farmer would actually receive, net of transport, shrinkage and the chance that a market is silent, and an adaptive conformal rule turned predictions into floors. Evaluation used a nested walk-forward protocol with Holm-corrected paired tests. Training on realised revenue cut regret by 87 to 157 rupees per quintal in all four crops, while a decision-focused loss added no reliable gain. Floors set to be breached on 10 percent of days were breached on 6.3 to 10.5 percent. The advice beat the highest price seen yesterday by 170 rupees per quintal for tomato and 202 for ginger, but not reliably for potato. An audit showed that the 27 to 62 percent skill of decompose-then-forecast hybrids falls to between minus 15 and minus 1 percent once the decomposition is causal. Grounded language-model briefs, five inspectable features (cost breakdown, what-if, track record, silence tax and crowd simulator) and a Telegram morning digest in English, Kannada and Hindi complete the system. A cross-commodity spillover engine and a council of forecasting agents did not meet the evidence bar and are reported as such. The limits are one region, one data source and no field trial.
